\rhead{RB202: Sensors, Actuators \& Signal Processing}
\phantomsection 
\section*{Sensors, Actuators \& Signal Processing (RB202)}
\label{major:RB_2}

\noindent \small{\hyperref[main:scheme]{$\leftarrow$ Back to Scheme}} \vspace{0.5cm}

\noindent \textbf{Credits:} 3 (2-0-2)

\subsection*{Theory}
\noindent\textbf{Unit 1} \\
\textit{Sensor Fundamentals and Classification:} Sensor principles and transduction mechanisms (resistive, capacitive, inductive, piezoelectric, optical), Sensor specifications (sensitivity, resolution, accuracy, hysteresis, linearity), Active vs. passive sensors, Analog vs. digital sensors, Environmental factors affecting sensor performance (temperature drift, noise, EMI), Sensor fusion concepts and Kalman filtering introduction.

\vspace{0.3cm}

\noindent\textbf{Unit 2} \\
\textit{Physical Sensors for Robotics:} Position and displacement sensors (potentiometers, LVDT, optical encoders, resolvers), Proximity and range sensors (ultrasonic, infrared, laser time-of-flight, LIDAR), Force/torque sensors (strain gauges, piezoelectric), Temperature sensors (thermocouples, RTDs, thermistors, infrared), Accelerometers and gyroscopes (MEMS principles, frequency response).

\vspace{0.3cm}

\noindent\textbf{Unit 3} \\
\textit{Signal Conditioning and Interface Circuits:} Analog signal conditioning (amplification, filtering, offset compensation), Instrumentation amplifiers and bridge circuits, Anti-aliasing filters and sampling theory, ADC/DAC fundamentals and selection criteria, Digital signal preprocessing (decimation, FIR/IIR filtering), Excitation circuits for active sensors (constant current/voltage sources).

\vspace{0.3cm}

\noindent\textbf{Unit 4} \\
\textit{Actuators and Drive Electronics:} DC motors (characteristics, torque-speed curves, commutation), Stepper motors (full/half step, microstepping), Servo motors (position feedback loops), Solenoids and voice coil actuators, Pneumatic/hydraulic actuators, Piezoelectric actuators, H-bridge drivers, PWM techniques and motor control algorithms.

\vspace{0.3cm}

\noindent\textbf{Unit 5} \\
\textit{Digital Signal Processing for Sensor Data:} Time-domain analysis (signal averaging, peak detection), Frequency-domain analysis (FFT fundamentals, spectral leakage), Digital filters (FIR/IIR design, stability), Noise reduction techniques (matched filtering, wavelet denoising), Sensor calibration and characterization, Real-time signal processing constraints and embedded implementation considerations.

\newpage

\subsection*{Practical}
\noindent\textbf{Lab 1: Sensor Characterization} \\
Focuses on measuring and documenting sensor specifications (sensitivity, linearity, hysteresis) using controlled test setups and data acquisition systems.

\vspace{0.2cm}

\noindent\textbf{Lab 2: Wheatstone Bridge and Strain Gauges} \\
Focuses on implementing instrumentation amplifier circuits for strain gauge measurements and characterizing force/load relationships.

\vspace{0.2cm}

\noindent\textbf{Lab 3: IMU Sensor Fusion} \\
Focuses on interfacing accelerometers and gyroscopes, implementing complementary and Kalman filters for attitude estimation.

\vspace{0.2cm}

\noindent\textbf{Lab 4: Ultrasonic and IR Ranging} \\
Focuses on comparing time-of-flight vs. triangulation ranging methods, characterizing accuracy vs. distance profiles.

\vspace{0.2cm}

\noindent\textbf{Lab 5: Anti-Aliasing Filter Design} \\
Focuses on designing and testing analog anti-aliasing filters, verifying Nyquist compliance with oscilloscope measurements.

\vspace{0.2cm}

\noindent\textbf{Lab 6: ADC Performance Analysis} \\
Focuses on characterizing ADC resolution, sampling rate, and ENOB through signal reconstruction and SNR measurements.

\vspace{0.2cm}

\noindent\textbf{Lab 7: DC Motor Characterization} \\
Focuses on measuring motor torque-speed curves, implementing open-loop speed control using PWM duty cycle variation.

\vspace{0.2cm}

\noindent\textbf{Lab 8: Servo Motor Position Control} \\
Focuses on implementing PID position control loops for servo motors with setpoint tracking and disturbance rejection analysis.

\vspace{0.2cm}

\noindent\textbf{Lab 9: Digital Filter Implementation} \\
Focuses on designing and implementing FIR/IIR filters for sensor noise reduction, comparing frequency responses.

\vspace{0.2cm}

\noindent\textbf{Lab 10: Sensor-Actuator Integration} \\
Focuses on building complete closed-loop systems combining multiple sensors with actuator control, implementing sensor fusion and real-time processing.

\vspace{0.2cm}

\newpage
\rhead{Curriculum Schema}
